\documentclass[11pt]{article}
\usepackage[T1]{fontenc}
\usepackage{lmodern,amsmath,amssymb,amsthm,mathtools,booktabs,microtype}
\usepackage[letterpaper,margin=1in,top=.9in,bottom=.92in]{geometry}
\usepackage{xcolor,hyperref}
\definecolor{ink}{HTML}{14394B}
\hypersetup{colorlinks=true,linkcolor=ink,citecolor=ink,urlcolor=ink,
pdftitle={Zeros and poles of translated Dirichlet L-functions in triple-logarithmic windows},pdfauthor={OpenAI}}
\usepackage{fancyhdr}
\pagestyle{fancy}\fancyhf{}\fancyhead[L]{\small\textcolor{ink}{Dirichlet zero divisors}}\fancyhead[R]{\small Local independence}\fancyfoot[C]{\thepage}
\setlength{\headheight}{14pt}\setlength{\emergencystretch}{2em}
\newtheorem{theorem}{Theorem}
\newtheorem{lemma}[theorem]{Lemma}
\newtheorem{corollary}[theorem]{Corollary}
\DeclareMathOperator{\sech}{sech}\DeclareMathOperator{\arsinh}{arsinh}
\DeclareMathOperator{\conv}{conv}\DeclareMathOperator{\ord}{ord}
\newcommand{\R}{\mathbb R}\newcommand{\C}{\mathbb C}\newcommand{\Z}{\mathbb Z}
\newcommand{\bs}{\boldsymbol}\newcommand{\one}{\mathbf 1}
\newcommand{\dd}{\,\mathrm d}
\newcommand{\lthree}{\log_3 T}
\newcommand{\wh}{\widehat}
\title{\vspace{-1.2em}\color{ink}\bfseries Zeros and poles of translated Dirichlet $L$-functions\\ in triple-logarithmic windows}
\author{Research note prepared for Vladimir Reshetnikov}
\date{26 November 2026}
\begin{document}
\maketitle
\begin{abstract}
Fix $B>0$ and a primitive Dirichlet $L$-function. In every sufficiently
high window of radius $\pi^{-1}\log\log\log T+O(1)$, the quotient
$L(s+ih)/L(s)$ has both a zero and a pole for every $0<|h|\le B$.
One window works for all steps simultaneously, however close to zero.
For steps below any fixed fraction of the mean zero spacing, the radius
shrinks to $O(1/\log\log T)$. Both results extend to finite differences
of every fixed order. The first uses two prime frequencies and uniform
control of the truncated explicit formula after division by a power of
the step. The second follows from a zero count and a divided-difference
argument. Dirichlet GRH supplies the critical-line input.
\end{abstract}

\section{The local statement}
Fix distinct primitive Dirichlet characters $\chi_1,\ldots,\chi_d$ of
conductors $q_1,\ldots,q_d$ and put $Q=\operatorname{lcm}(q_1,\ldots,q_d)$. The primitive character of conductor one gives
$\zeta$. For a nontrivial zero $\rho$ of at least one member of the family, put
\[
 m(\rho)=(\ord_\rho L(s,\chi_1),\ldots,\ord_\rho L(s,\chi_d)).
\]
A missing zero contributes zero to its coordinate. All sums indexed by the
union of the zero sets use each distinct point once. Under Dirichlet GRH write
$m(\gamma)=m(1/2+i\gamma)$. Write $\lthree=\log\log\log T$, using natural
logarithms, and $x_+=\max(x,0)$. Constants may depend on this fixed family.

\begin{theorem}[Uniformity as the shifts coalesce]\label{th:diff}
Fix a primitive character $\chi$, an integer $r\ge1$, and $B>0$.
For $0<|h|\le B$ form the meromorphic function
\[
 F_{r,h}(s)=\prod_{k=0}^r L(s+ikh,\chi)^{(-1)^{r-k}\binom rk}.
\]
There exist $C,T_0$ independent of $h$ such that, for every $T\ge T_0$,
$F_{r,h}$ has both a zero and a pole on the critical line in
$|\Im s-T|\le\pi^{-1}\lthree+C$. If $N_+,N_-$ count zeros and poles
in that window after cancellation, with multiplicity, then
$N_\pm\ge c|h|^r$. For any fixed $b>1/\pi$, in
$|\Im s-T|\le b\lthree$ one has
\[
 N_\pm\ \gg_{\chi,r,B,b}\ \max\{1,|h|^r\lthree\}.
\]
These statements use Dirichlet GRH. One threshold $T_0$ works for
all nonzero steps in the indicated range.
\end{theorem}

For example $\zeta(s+ih)/\zeta(s)$ has both a zero and a pole for every
$0<|h|\le1$. For $r=2$ the quotient is
$\zeta(s+2ih)\zeta(s)/\zeta(s+ih)^2$. At $h=0$ all factors cancel.
The proof estimates the explicit formula after division by $h^r$.
Theorem~\ref{th:micro} below gives a shrinking radius $C/\log\log T$
when $|h|\le\kappa/\log T$ for any fixed $\kappa<2\pi$. Corollary~\ref{cor:nonres} extends the shrinking-window result to bounded multiples of the mean spacing separated from the positive integers.

\begin{theorem}[Both signs in every high window]\label{th:sign}
Let $d\ge2$. There exist $C,c>0$ and $T_0$ such that for all $T\ge T_0$,
every real unit vector $u\perp\one$, and each $s\in\{-1,1\}$,
\begin{equation}
 \sum_{|\gamma-T|\le \pi^{-1}\lthree+C}
      \big(s\, u\cdot m(\gamma)\big)_+\ \ge c .
 \label{eq:sign}
\end{equation}
For every fixed $b>1/\pi$ and all sufficiently large $T$, one also has
\begin{equation}
 \sum_{|\gamma-T|\le b\lthree}
      \big(s\, u\cdot m(\gamma)\big)_+\ \gg_{b,\chi_1,\ldots,\chi_d} \lthree .
 \label{eq:mass}
\end{equation}
Both assertions use the critical-line theorem for Dirichlet $L$-functions.
\end{theorem}

In particular, for any nonzero $n\in\Z^d$ with $\sum_j n_j=0$, the
meromorphic function $\prod_j L(s,\chi_j)^{n_j}$ has both a zero and a pole
in every window in \eqref{eq:sign}. All cancellations at shared zeros have
already been made. On the windows of \eqref{eq:mass}, its numbers of zeros and
poles, counted with multiplicity, are each $\gg \|n\|_2\lthree$, uniformly in
$n$. Taking $n=(1,-1)$ gives the assertion for a quotient of two functions.
A larger zero order for one member permits a positive zero order for the other.

\begin{theorem}[Local full rank, independent of GRH]\label{th:rank}
There exist $C,T_0$ such that for $T\ge T_0$ the vectors
\[
 \big\{m(\rho):\ |\Im\rho-T|\le \pi^{-1}\lthree+C\big\}
\]
span $\C^d$. Here the index is the full complex point $\rho$. In particular any
nonzero complex linear combination of $L'/L(s,\chi_j)$ has a pole in the
indicated height range.
\end{theorem}
Distinct primitive characters are essential. Inducing the same primitive
character to different moduli gives the same zeros in the open critical strip.
The constants and the starting height in these theorems are family-dependent.
The coefficient $1/\pi$ is achieved by the proof. Its optimality for zero
divisors is open.

\section{Sources and comparison}
The public theorem \texttt{LFunction\_eq\_zero\_re} in the readable
\texttt{grh-lean} project places every open-critical-strip zero of every Dirichlet
$L$-function on $\Re s=1/2$. Its companion theorem excludes zeros in the right
half of the strip directly. In its reduction a one-sided prime-power estimate
is obtained by combining short integral forms, a spanning theorem and a
determinant lower bound. We use the exported zero-location statement and
work in the other direction of the explicit formula, choosing prime tests
that force a signed zero discrepancy. The repository and its independent
verification records distinguish these Dirichlet results from hypotheses for
general automorphic $L$-functions \cite{lean,checks}. The linked overview
describes the explicit-formula connection and the status of the formal proof
\cite{wiki}.

Srinivas proved for distinct members $F,G$ of the Selberg class, with
$\deg F\ge \deg G$, that every high interval $|\Im\rho-T|\le C\log\log T$
contains a point with $m_F(\rho)>m_G(\rho)$ \cite{srinivas}. This applies
in both directions for a pair of degree one. His proof already uses a rapidly
decaying Ramachandra kernel. Theorem~\ref{th:sign} shortens the radius in the
Dirichlet case and handles all real balanced directions of a fixed family.
Kaczorowski and Perelli's global multiplicity-basis theorem \cite{kp}, recorded as
Lemma 7.1 in \cite{perelli}, gives infinitely many independent tuples.
Theorem~\ref{th:rank} prescribes the location of a complete tuple at every high
center. Global positive-proportion results, conductor-uniform comparisons,
and unshared-zero theorems address different quantitative questions
\cite{bp,banks,bl}.

The ProveIt directory \cite{proveit} catalogs 61 initial and 42 subsequent
packets, alongside staged material. Its packet on balancing prime errors
proves global spanning of multiplicity vectors. A Gaussian prime-frequency
packet gives a $C\log\log T$ reconstruction window. Two closer predecessors
already reach the triple logarithm. \emph{Reading arithmetic weights from
poles of a prime series} \cite{pole} gives local full rank for every leading
constant $b>1/\pi$, and both signs for a pair of the same primitive conductor.
\emph{Identifying Dirichlet characters from a short window of zeros}
\cite{euler} recovers individual prime coefficients from complex zeros at
the endpoint $1/\pi$, with explicit conductor and accuracy dependence.
Its linear measurements also yield endpoint rank for a fixed family by
choosing finitely many distinguishing primes (the common gamma correction
tends to zero). It gives a matching-cost bound depending on conductor.

A particularly close predecessor is the section on multiplicity differences
in \cite{short}. For fixed unshifted primitive characters it gives both
signs for every balanced real direction, including unequal conductors, and
mass $\gg H\log\log H$ on intervals of length $H\ge(2/\pi+\epsilon)\lthree$.
This is stronger in mass than \eqref{eq:mass} for those unshifted families.
Our fixed prime tests give the additive-constant endpoint in \eqref{eq:sign}
and permit uniform estimates after translating the functions.

The proposed contributions focus on moving shifts, especially the bounds
through coalescence in Theorems~\ref{th:diff} and \ref{th:micro}.
Ramachandra and Sankaranarayanan and subsequent work already treated zeros
or poles of quotients of translated zeta products \cite{trans,brss}.
The specific result here uses a single window and starting height for all
nonzero steps up to $B$, and locates both signs in shrinking windows for
steps below the mean spacing. We give an independent proof of the supporting
rank statement and a binomial variant of a critical-strip cutoff. The
previously known unshifted scale, rank, and sign conclusions are credited
above. Publication priority for the new statements remains to be established.

\section{A test function localized across the critical strip}

Use $\wh h(\xi)=\int_\R h(t)e^{-it\xi}\dd t$. We recall the
Guinand--Weil explicit formula for a primitive character $\chi$ of conductor
$q$ and parity $e\in\{0,1\}$. Set $z_\rho=(\rho-1/2)/i$ and
\[
 G_\chi(t)=\log(q/\pi)+\Re\frac{\Gamma'}{\Gamma}
                 \left(\frac{1/2+e+it}{2}\right).
\]
For tests holomorphic and rapidly enough decreasing in a strip wider than
$|\Im z|\le1/2$, the convention used here gives
\begin{align}
 2\pi\sum_{\rho\text{ of }\chi}h(z_\rho)
   &=2\pi\delta_{q=1}\{h(i/2)+h(-i/2)\}
      +\int_\R h(t)G_\chi(t)\dd t \notag\\[-2pt]
   &\quad -\sum_{n\ge2}\frac{\Lambda(n)}{\sqrt n}
   \big\{\chi(n)\wh h(\log n)+\overline{\chi(n)}\wh h(-\log n)\big\}.
   \label{eq:weil}
\end{align}
The sum on the left here alone counts zeros with multiplicity. This is the
classical explicit formula for the completed function, applied to a complex
test by linearity \cite{rs,dav}. For the tests below, one can pass from compactly supported smooth Fourier transforms by smooth truncation at frequency $N$. For each fixed test its Fourier transform and its first two derivatives decay as $O(e^{-2|\xi|/3})$, by shifting its defining integral to height $\pm2/3$. Two integrations by parts show convergence of the inverse transforms in the norm $\sup_{|y|\le1/2,x}(1+x^2)|h(x+iy)|$. The unit-interval zero count and absolute convergence on the prime side justify the limit. It follows also by shifting the contour of
the completed logarithmic derivative, using the functional equation on the
left boundary. The simple poles at $0,1$ for $\zeta$ give the displayed
pair. The classical unit-interval zero bound, with multiplicity, is
$O_q(\log(2+|t|))$ on $[t,t+1]$ \cite{dav}.

\begin{lemma}[Critical-strip cutoff]\label{lem:kernel}
Fix $W\ge1$. There are even probability densities $K=K_{M,W}$ for integers
$M\ge2$ with the following estimates, uniform in $M$:
\begin{align}
 |K(x+iy)|&\ll_W e^{-2|x|} &&(|y|\le1/2),\label{eq:k1}\\
 |K(x+iy)|&\ll_W 2^{-M/2}e^{-2|x|}
       &&(|y|\le1/2,\ |x|\ge W+\pi^{-1}\arsinh M),\label{eq:k2}\\
 |\wh K(\xi)|&\ll e^{-|\xi|/2}\min\{1,(4/(W|\xi|))^4\}.
       \label{eq:khat}
\end{align}
Each $K$ is admissible in \eqref{eq:weil}, is nonnegative on the real axis,
and has $\wh K(0)=1$.
\end{lemma}
\begin{proof}
Normalize
\[
 k_M(z)=c_M\sech^2(z)\big(1+M^{-1}\cosh(\pi z)\big)^{-M}
\]
to have real integral one. The $c_M$ are bounded above and below uniformly,
by comparison with $\sech^2 x$ and by integration over $[-1,1]$ using
$(1+y/M)^M\le e^y$ for $y\ge0$. For $z=x+iy$ in the closed critical strip,
$w=\cosh(\pi z)$ has $\Re w\ge0$ and
$|w|^2=\sinh^2(\pi x)+\cos^2(\pi y)$. Thus $|1+w/M|\ge1$, and
$|1+w/M|\ge\sqrt2$ when $|w|\ge M$.
Also $|\sech^2(x+iy)|\ll e^{-2|x|}$ there.
The nearest possible poles of the binomial factor have $|\Im z|=1$:
if $\cosh\pi z=-M$, its imaginary part forces $\sin\pi y=0$,
since $x=0$ is impossible for $M>1$. The $\sech$ poles are farther away.
For each fixed $M$ the function therefore admits a wider strip, for example
$|\Im z|\le3/4$, with exponential decay as $|x|\to\infty$.
Shifting the Fourier integral to $\Im z=-\operatorname{sgn}(\xi)/2$
gives $|\wh k_M(\xi)|\ll e^{-|\xi|/2}$ with an absolute constant.

Convolve $k_M$ with four uniform probability densities on $[-W/4,W/4]$.
The resulting compact average of translates preserves holomorphy and
gives \eqref{eq:k1}--\eqref{eq:k2}. Its Fourier multiplier is
$\big(\sin(W\xi/4)/(W\xi/4)\big)^4$, proving \eqref{eq:khat}.
\end{proof}

For all subsequent estimates at center $U$, with $T/2\le U\le2T$, take
\[
 M=\lceil 10\log\log T\rceil,\qquad
 H=W+\pi^{-1}\arsinh M=\pi^{-1}\lthree+O_W(1).
\]
Multiplying
$K(z-U)$ by any fixed finite sum of $e^{i\lambda z}$ with bounded
coefficients costs a fixed factor on $|\Im z|\le1/2$. By
\eqref{eq:k2} and the unit-interval zero bound, the tail
$|\Im\rho-U|>H$ of the resulting zero sum is
\begin{equation}
  O_W\!\left(2^{-M/2}\log(2T)\right)=o(1).
 \label{eq:tail}
\end{equation}
Indeed the sum of $e^{-2|\Im\rho-U|}$ is $O_q(\log(2T))$ by
unit-interval summation. The pole terms are $O_W(e^{-2U})$.
Stirling's formula on $|t-U|\le U/2$, followed by \eqref{eq:k1} on
its complement, gives, for all these tests,
\begin{equation}
 \int h(t)G_\chi(t)\dd t =
     \log(qU/(2\pi))\int h(t)\dd t+O_W(T^{-1}).
 \label{eq:gamma}
\end{equation}
Here the pointwise error in the central range is
$O((1+|t-U|)/U)$ and the first absolute moment of $K(t-U)$ is uniformly
bounded for fixed $W$. In estimates with prescribed frequencies we allow
constants also to depend on that finite list.

\section{Selecting prime frequencies}

For $t\mapsto K(t-U)e^{i\omega t}$ the transform at $\xi$ is
$e^{iU(\omega-\xi)}\wh K(\xi-\omega)$. Consequently the prime sum in
\eqref{eq:weil} tends, as $W\to\infty$, to its terms with
$\omega=\log n$ or $\omega=-\log n$, with a uniform error $o_W(1)$
over $U$ and $M$. At an exact match the multiplier is $1$. To justify
discarding every other term simultaneously, use \eqref{eq:khat}. For fixed
$\omega$ and $n$ large, either sign has the summable majorant
\[
 C_\omega\frac{\log n}{n(1+\log n)^4}.
\]
For any fixed unmatched $n$ the fourth power in \eqref{eq:khat} makes its
contribution tend to zero uniformly. This argument works for any fixed finite
list of frequencies, including zero. At nonzero $\omega$, the real integral
of the test also tends to zero uniformly, since its modulus is
$|\wh K(\omega)|$. The order of choices below matters: we fix the primes,
then choose $W$, and only then take $T$ large.

\begin{proof}[Proof of Theorem~\ref{th:rank}]
The characters induced on the
units modulo $Q$ are distinct, so their value table has column rank $d$.
Choose $d$ independent rows, and, by Dirichlet's theorem, one prime in
each corresponding reduced residue class. Write $p_1,\ldots,p_d$ for these
primes. The matrix $B_{ij}=\chi_j(p_i)\log(p_i)/\sqrt{p_i}$ is invertible.
For $\lambda=\log p_i$ use the test
\[
 h_i(t)=K(t-T)\big(e^{i\lambda t}-e^{i\lambda T}\wh K(\lambda)\big).
\]
Its integral is exactly zero by evenness of $K$. Thus \eqref{eq:gamma}
is $O_W(1/T)$ despite the growing common gamma term. Prime-frequency
selection just proved gives the matrix estimate
\[
 2\pi\!\sum_{|\Im\rho-T|\le H}h_i(z_\rho)m_j(\rho)
   =-B_{ij}+o_W(1)+o_{T;W}(1).
\]
The unmodulated correction has a uniformly bounded prime sum and a
coefficient $o_W(1)$. The tail uses the strip estimate \eqref{eq:tail},
including zeros away from the center line. Choose $W$ and then $T_0$ to
make the error in operator norm less than half the smallest singular value
of $B$. The left-hand matrix has rank $d$. Its rows are in the complex
span of the indicated multiplicity vectors, as required. At a zero $\rho$,
the residue of $\sum_j v_j L'/L(s,\chi_j)$ is $v\cdot m(\rho)$;
the spanning statement gives the last assertion.
\end{proof}

The zero-location input enters precisely at the next step. It makes each
test of the form $K(t-U)|A(t)|^2$ nonnegative at every zero argument.
The common gamma contribution grows like $\log T$. The condition
$\sum_j u_j=0$ removes it, leaving a bounded conductor contribution that
can have either sign.

\begin{lemma}[A finite set of distinguishing primes]\label{lem:primes}
For any $D>0$ there is a finite set $P$ of primes coprime to $Q$ such that
\begin{equation}
 \sum_{p\in P}\frac{(\log p)^2}{p}\left|\sum_j u_j\chi_j(p)\right|^2
       > D^2\|u\|_2^2 \qquad (u\in\C^d\setminus\{0\}).
 \label{eq:frame}
\end{equation}
\end{lemma}
\begin{proof}
In every reduced class $a$ modulo $Q$, the sum of $1/p$ over its primes
diverges, by the prime number theorem in arithmetic progressions
(or Dirichlet's reciprocal-prime theorem). Choose an initial prime list
so that the sum of $(\log p)^2/p$ exceeds $D^2/\varphi(Q)$ in
every class. Summing over classes and using character orthogonality proves
\eqref{eq:frame}.
\end{proof}

\section{Positive tests for both directions}

For $u\in\R^d$ of norm one with $u\cdot\one=0$ put
\[
 c_u=\sum_j u_j\log q_j,\qquad
 b_p(u)=\frac{\log p}{\sqrt p}\sum_j u_j\chi_j(p).
\]
Fix $P$ by Lemma~\ref{lem:primes} so that
$\|b(u)\|_2>\| (\log q_j)_j\|_2+3$ for every such $u$.
For $a=(a_0,(a_p)_{p\in P})$ of complex Euclidean norm one, set
\[
 A_a(t)=a_0+\sum_{p\in P}a_p e^{it\log p},\qquad
 h_{a,U}(t)=K(t-U)|A_a(t)|^2 \quad(t\in\R).
\]
For use in \eqref{eq:weil} the holomorphic extension of the squared modulus
is $A_a(z)\overline{A_a(\bar z)}$. Its absolute value in the strip is
bounded uniformly in $a$, for the fixed $P$.

Expanding the square gives frequency zero with coefficient one. The
frequencies $\log p$ and $-\log p$ come from the constant term paired
with $p$. The remaining frequencies are $\log(p/r)$ for distinct primes
$p,r$. Neither $p/r$ nor $r/p$ is an integer, so none of these remaining
frequencies equals a signed prime-power logarithm. In particular
\[
 \int h_{a,U}(t)\dd t=1+o_W(1).
\]
Combine \eqref{eq:weil} over $j$ with coefficients $u_j$. The gamma
contribution is $c_u\int h_{a,U}+O_W(T^{-1})$ by \eqref{eq:gamma}.
The frequency selection in the preceding section and \eqref{eq:tail} give
\begin{align}
 2\pi\sum_{|\gamma-U|\le H} h_{a,U}(\gamma)\, u\cdot m(\gamma)
 = c_u-2\Re\!\left(\overline{a_0}\sum_{p\in P}a_p b_p(u)\right)
       +o_W(1)+o_{T;W}(1).\label{eq:star}
\end{align}
All errors are uniform in $T/2\le U\le 2T$, in $a$ and in $u$.
This follows from the finite expansion, $\|a\|=\|u\|=1$, and the
uniformity of the frequency and strip bounds. In particular the
$\log T$ contribution cancels \emph{before} any frequency leakage is bounded.

Take $a_0=1/\sqrt2$ and $a_p=\pm\overline{b_p(u)}/(\sqrt2\|b(u)\|_2)$.
The two main terms in \eqref{eq:star} are $c_u\mp\|b(u)\|_2$.
Each has its asserted sign with a margin greater than three. Fix $W$ to
make the first error less than one, then $T_0$ to make the second less than
one. Since $|A_a(t)|^2\le 1+|P|$ and $K\ge0$, we have proved a weighted
form of both inequalities, with $c_0=1/(2\pi(1+|P|))$:
\begin{equation}
 \sum_{|\gamma-U|\le H}K(\gamma-U)
       (s\,u\cdot m(\gamma))_+\ge c_0,\qquad s=\pm1 .
 \label{eq:weighted}
\end{equation}
The uniform supremum bound on $K$ and $U=T$ prove \eqref{eq:sign}.
For \eqref{eq:mass} put $R=b\lthree$. In \eqref{eq:weighted}, $M,W,H$
are fixed at their values for $T$. Integrate over $|U-T|\le R-H$.
All contributing ordinates satisfy $|\gamma-T|\le R$, and a given
ordinate receives integrated weight at most $\int_\R K=1$. Consequently
\[
 \sum_{|\gamma-T|\le R}(s\,u\cdot m(\gamma))_+
 \ \ge\ 2c_0(R-H) \ \gg_{b,\chi_1,\ldots,\chi_d}\lthree .
\]
This completes the proof of Theorem~\ref{th:sign}. The enlargement from
$1/\pi$ to $b$ in the mass estimate provides the integration interval.

\section{Two consequences for local zero data}

\begin{corollary}[Exact balancing]\label{cor:conv}
In each window of Theorem~\ref{th:sign}, let
$v_\gamma=m(\gamma)/\sum_jm_j(\gamma)$. In the affine hyperplane
$\sum_j x_j=1$, the convex hull of the $v_\gamma$ contains a Euclidean
ball of radius $c_1/\log T$ about $\one/d$.
There are at most $d$ ordinates and positive integer weights on them
whose weighted sum of multiplicity vectors has all coordinates equal
to the same positive integer.
\end{corollary}
\begin{proof}
The total weight $\sum_{|\gamma-T|\le H}K(\gamma-T)\sum_jm_j(\gamma)$
is $O(\log T)$ by \eqref{eq:k1} and the unit-interval bound.
Equation \eqref{eq:weighted} implies
$\max_\gamma u\cdot v_\gamma\gg1/\log T$ for every unit
$u\perp\one$ in this window. Separation by a hyperplane gives the ball.
Carath\'eodory's theorem yields a convex combination equal to $\one/d$
on at most $d$ affinely independent points. Its barycentric coordinates
are rational by linear algebra over $\mathbb Q$; discard zero coefficients.
Divide each coefficient by $\sum_jm_j(\gamma)$ and clear denominators
to obtain the integer weights.
\end{proof}

\begin{corollary}[A cost for moving or deleting zeros]\label{cor:transport}
For two fixed distinct primitive characters, form the multisets of both
zero lists in $I=[T-b\lthree,T+b\lthree]$, $b>1/\pi$. For any partial
one-to-one matching of the two lists, charge $1$ for each unmatched
occurrence and $\min(1,|\gamma-\gamma'|)$ for each matched pair. Every
matching has cost $\gg_{b,\chi_1,\chi_2}\lthree$.
\end{corollary}
\begin{proof}
Use $u=(1,-1)/\sqrt2$ and choose the positive test in \eqref{eq:star}.
Its coefficients depend only on $u$ and $P$. For $H$ as in the proof let
$\eta_H(x)$ be one on $[-H,H]$, linear down to zero on
$H\le |x|\le H+1$, and zero outside. In \eqref{eq:star} we may use this
taper in place of the indicator, by \eqref{eq:tail}. For
$J=[T-R+H+1,T+R-H-1]$ with $R=b\lthree$, form
\[
 f(t)=|A_a(t)|^2\int_J K(t-U)\eta_H(t-U)\dd U .
\]
It is nonnegative and supported in $I$. Its supremum and Lipschitz
constant are bounded independently of $T$: the inner factor is bounded
by one and is a difference of two values of an antiderivative of
$K\eta_H$, with derivative bounded by $2\|K\|_\infty$.
The finite polynomial and its derivative are bounded on the real line.
Integrating the signed estimate over $J$ gives
$\sum f(\gamma_1)-\sum f(\gamma_2)\gg |J|\gg\lthree$.
Any matched pair contributes at most
$\max(2\|f\|_\infty,\operatorname{Lip}f)\min(1,|\gamma_1-\gamma_2|)$
in absolute value. An unmatched zero contributes at most $\|f\|_\infty$.
Summation proves the claim.
\end{proof}

\section{Moving ordinates and coalescing shifts}\label{sec:moving}

Consider $L(s+i\tau_j,\chi_j)$ for real shifts $\tau=(\tau_1,\ldots,\tau_d)$
in a fixed nonempty compact set $E$. Allow repeated primitive characters
and assume that $(\chi_j,\tau_j)\ne(\chi_k,\tau_k)$ throughout $E$ for
$j\ne k$. The multiplicity at $1/2+it$ is $m_{\chi_j}(t+\tau_j)$.

\begin{theorem}[Uniform translated families]\label{th:moving}
Theorems~\ref{th:sign} and \ref{th:rank} and Corollary~\ref{cor:conv}
hold for these shifted functions uniformly in $\tau\in E$. For a pair,
Corollary~\ref{cor:transport} also holds uniformly. Its lower bound on
matching cost persists after taking the infimum over all allowed translations.
\end{theorem}
\begin{proof}
Testing the original list with $h(t-\tau_j)$ replaces its prime coefficient
by $\chi_j(n)n^{-i\tau_j}$ and its gamma density by $G_{\chi_j}(t+\tau_j)$.
The pole terms and all error bounds are uniform for bounded shifts.
It suffices to choose the primes in \eqref{eq:frame} uniformly.
For fixed $\tau$, the matrix
\[
 G_X(j,k)=\sum_{p\le X,\,p\nmid Q}\frac{(\log p)^2}{p}
 \chi_j(p)\overline{\chi_k(p)}p^{-i(\tau_j-\tau_k)}
\]
satisfies $G_X/(\tfrac12(\log X)^2)\to I_d$. Indeed a nonprincipal
character has $\sum_{p\le x}\chi(p)\log p=o(x)$ by the prime number
theorem in progressions. Partial summation against
$x^{-1-iv}\log x$ gives $o((\log X)^2)$ for fixed $v$. For the principal
character the main term is
$\int_2^X(\log x)x^{-1-iv}\dd x$, which is
$\tfrac12(\log X)^2+O(1)$ for $v=0$ and $O_v(\log X)$ otherwise.
When the primitive characters coincide off the diagonal, their shifts differ.
For a given $D$, the condition $\lambda_{\min}(G_X)>D^2$ is open in
$\tau$ and persists on increasing $X$, since the matrix increases as a
quadratic form. Compactness supplies a single $X$ valid throughout $E$.
The rank proof now uses a rectangular prime matrix with a uniform lower
singular value. All subsequent arguments use finitely many frequencies with
uniformly bounded coefficients. The polynomial in the signed test can depend
on $\tau$ and $u$; its supremum, Lipschitz bound and signed margin are uniform.
\end{proof}

We now prove Theorem~\ref{th:diff}. Here the shifted characters approach
a collision, so the preceding compactness argument supplies no lower
singular value. For a test $f$ define
\[
 \mathcal D_{r,h}(f)=h^{-r}\sum_{k=0}^r d_k
             \sum_{\gamma\text{ of }\chi} f(\gamma-kh),\qquad
             d_k=(-1)^{r-k}\binom rk .
\]
The zeros in the inner sum are repeated according to multiplicity.
The signed measure before normalization has atom $\ord_{1/2+it}F_{r,h}$ at $t$.
For a smooth function $g$ the identity
\begin{equation}
 h^{-r}\sum_{k=0}^r d_k g(x+kh)
 =\int_{[0,1]^r}g^{(r)}(x+h(v_1+\cdots+v_r))\dd v
 \label{eq:fd}
\end{equation}
follows by $r$ applications of the fundamental theorem of calculus,
and controls the difference even when $|d_k/h^r|$ becomes arbitrarily large.

Choose distinct primes $p_1,p_2$ prime to $q_\chi$, and put
\[
 b_p(h)=\chi(p)\frac{\log p}{\sqrt p}
       \left(\frac{e^{-ih\log p}-1}{h}\right)^r .
\]
Extend continuously to $h=0$. These two functions have no common zero
on $[-B,B]$: simultaneous vanishing for $h\ne0$ would make
$\log p_1/\log p_2$ rational. At zero both limits are nonzero. Thus
$\max_{p=p_1,p_2}|b_p(h)|\ge\delta$ for some $\delta>0$ throughout the interval.
For example, the certificate in the archive gives
$\sum_{p=2,3}|b_p(h)|^2>1/2$ for $\chi=1$, $r=2$, $|h|\le1$.

In Lemma~\ref{lem:kernel} use $J=r+4$ copies of the uniform probability
density on $[-W/J,W/J]$. The support radius of the compact average is still
$W$ and \eqref{eq:khat} has power $J$ and numerator $J$. Bounds
\eqref{eq:k1}--\eqref{eq:k2} hold also for $K^{(j)}$, $0\le j\le r$,
with constants depending on $r,W$: put one derivative on each of $j$
uniform factors. Its distributional derivative is a signed measure on
two endpoints, so the translate average of $k_M$ has uniformly
bounded total variation for fixed $r,W$.

Choose $p$ realizing the preceding maximum and use the two positive tests
\[
 f(t)=K(t-U)P(t),\qquad
 P(t)=1\pm\Re(\alpha e^{it\log p}),\qquad
 \alpha=\overline{b_p(h)}/|b_p(h)|.
\]
Extend the trigonometric polynomial by its two exponentials off the real line.
In the explicit formula for $\mathcal D_{r,h}(f)$ the prime coefficient
$\chi(n)$ is multiplied by $((n^{-ih}-1)/h)^r$, of magnitude at most
$(\log n)^r$. The large-$n$ majorant for the off-frequency contributions is
$O_p((\log n)^{r+1}/(n(1+\log n)^J))$, which is summable.
Thus the prime peak and its conjugate give
$\mp |b_p(h)|+o_W(1)$ in $2\pi\mathcal D_{r,h}(f)$, uniformly in the step.
On the gamma side the density is $h^{-r}\sum_k d_k G_\chi(t+kh)$.
Equation \eqref{eq:fd} bounds it by $O_{r,B,\chi}(T^{-r})$ for
$|t-U|\le U/2$, and by a constant globally: the differentiated Stirling
expansion gives $G_\chi^{(r)}(t)=O_{r,\chi}((1+|t|)^{-r})$ for $r\ge1$.
Equation \eqref{eq:k1} then gives integrated error $o_{T;W}(1)$.
For the zeta pole terms use \eqref{eq:fd} on horizontal translates of
$f$ at $\pm i/2$ and the derivative bounds above, for total
$O_{r,B,W}(e^{-2U})$.

Choose a smooth $0\le\eta_H\le1$, equal to one on $[-H,H]$, zero outside
$[-H-1,H+1]$, with derivatives through order $r$ bounded independently
of $H\ge1$. Translate a fixed smooth cutoff on each side.
Put $g(t)=(1-\eta_H(t-U))f(t)$. On the real axis every derivative
through order $r$ of $g$ is $O_{r,W}(2^{-M/2}e^{-2|t-U|})$:
it vanishes inside $|t-U|<H$ and elsewhere obeys \eqref{eq:k2} and
its derivative versions. Apply \eqref{eq:fd} before summing over the zeros.
The bounded shift $|kh|\le rB$ and the unit-interval zero count give
\[
 |\mathcal D_{r,h}(g)|\ll_{r,B,W}2^{-M/2}
          \sum_\gamma e^{-2|\gamma-U|}=o(1).
\]
Truncation has introduced no factor $|h|^{-r}$. Fix $W$ to make the prime
leakage smaller than $\delta/4$ and take $T$ large for the other errors.
The two tapered nonnegative tests $\eta_H f$ have normalized signed
readings of opposite signs and magnitude bounded below. Multiplication by
$h^r$ may interchange the signs. Since $0\le P\le2$, for $s=\pm1$ this gives
\[
 \sum_{|t-U|\le H+1}K(t-U)(s\,\ord_{1/2+it}F_{r,h})_+
               \gg_{\chi,r,B}|h|^r .
\]
The uniform supremum bound gives both signs in the endpoint window.
Integration over $U$ as in \eqref{eq:mass} gives the larger-window bound
$\gg |h|^r\lthree$. Each positive count is an integer and hence at least
one. This proves Theorem~\ref{th:diff}.

\section{Steps shorter than the mean zero spacing}

There is a second mechanism when the step is smaller than the local
mean spacing $2\pi/\log T$. In this range the observation window shrinks
as the height grows. This result uses a zero count and ordinary divided
differences, and is independent of the prime-frequency proof.

\begin{theorem}[Microscopic steps]\label{th:micro}
Fix $r\ge1$, a primitive $\chi$, and $0<\kappa<2\pi$. For any
\[
 C>\frac{\pi}{2}+\frac{\pi\kappa}{2\pi-\kappa}
\]
and every sufficiently large $T$, the function $F_{r,h}$ has both a
zero and a pole in $|\Im s-T|\le C/\log\log T$, simultaneously for all
$0<|h|\le\kappa/\log T$. In particular every interval $[T,T+1]$ contains
$\gg_{\kappa,\chi,r}\log\log T$ zeros and as many poles, counted after
cancellation, uniformly for those steps.
\end{theorem}
\begin{proof}
Put $L=\log T$ and $\ell=\log\log T$. The degree-one argument estimate of
Carneiro, Chandee and Milinovich under Dirichlet GRH gives
$|S_\chi(t)|\le(1/4+o(1))L/\ell$ uniformly for $|t-T|=O(1)$
\cite{ccm}. The one-sided Riemann--von Mangoldt formula, or the argument
principle on the critical line with the gamma factor, consequently gives
for fixed $a>0$ and $|x-T|\le2$
\[
 \left|N_\chi((x,x+a/\ell))-\frac{a}{2\pi}\frac{L}{\ell}\right|
       \le(1/2+o(1))\frac{L}{\ell}.
\]
Indeed the smooth phase increment is $(a/\ell)(L+O_\chi(1))/(2\pi)$
and the two argument terms contribute at most the stated error.
The bound applies to open or closed endpoints by one-sided limits.
Choose a central interval of length $a/\ell$ with midpoint $T$ and
$a>\pi$. It contains a zero of ordinate $\gamma_0$ by the lower bound.

For $A>\pi\kappa/(2\pi-\kappa)$ consider each of the two adjacent
intervals of length $A/\ell$ starting from $\gamma_0$. On the grid
$\gamma_0+nh$, each side has at least $A L/(\kappa\ell)-O(1)$ sites.
The total number of zeros on that side, even counted with multiplicity,
is at most $(A/(2\pi)+1/2+o(1))L/\ell$. The first coefficient is
strictly larger than $A/(2\pi)+1/2$ by the choice of $A$. Thus each side
has at least $r$ empty grid sites for sufficiently large $T$, uniformly in
the step. Write $v(n)=m_\chi(\gamma_0+nh)$. Then $v(0)>0$.

Take $r$ negative empty indices and the index $0$. The divided difference
of order $r$ of $v$ on these $r+1$ nodes is positive: in the Lagrange
formula only $v(0)$ survives and its denominator is positive. On $r-1$
negative empty indices, $0$, and one positive empty index the result
is negative, by the same formula. We use signs in the integer index
coordinate, so this works for either sign of $h$.

Every divided difference on ordered integer nodes $x_0<\cdots<x_r$
is a convex combination of divided differences on $r+1$ consecutive
sites within their span. For a missing interior integer $y$, interpolate
between the divided differences on the node lists
$(x_0,\ldots,x_{r-1},y)$ and $(y,x_1,\ldots,x_r)$ with weights
$(y-x_0)/(x_r-x_0)$ and $(x_r-y)/(x_r-x_0)$ respectively.
The Lagrange formula gives the identity, independently of the ordering
inside each list. Both spans decrease, so induction ends at consecutive
nodes. On consecutive sites the divided difference is $\Delta^r v(n)/r!$,
where $\Delta v(n)=v(n+1)-v(n)$. The positive and negative readings
therefore force consecutive-grid differences of both signs. These are
exactly the orders of $F_{r,h}$ at the corresponding ordinates.

All those ordinates are within $(a/2+A)/\ell$ of $T$.
The strict hypothesis on $C$ permits choices with $a/2+A<C$.
For the unit-interval assertion use disjoint windows of this radius
inside $[T,T+1]$, with slightly larger fixed $\kappa'<2\pi$
to cover the variation of $\log T$ across the interval.
Each window contributes at least one zero and one pole.
\end{proof}

This proof uses $\kappa<2\pi$ only to produce more grid sites than zero
occurrences. A lattice with spacing exactly $2\pi/\log T$ is invariant
under that translation, so a density comparison alone could not give
the conclusion at equality. Sharpness for the actual zeros remains open.

\begin{corollary}[Nonintegral multiples of the mean spacing]\label{cor:nonres}
Let $K_0>0$ and $\delta>0$ be fixed. For
$\nu=|h|\log T/(2\pi)$ suppose $0<\nu\le K_0$ and
$\operatorname{dist}(\nu,\{1,2,\ldots\})\ge\delta$.
Both a zero and a pole of $F_{r,h}$ occur in
$|\Im s-T|\le C_{\chi,r,K_0,\delta}/\log\log T$, uniformly in $h$.
\end{corollary}
\begin{proof}
For $\nu\le1/2$ apply Theorem~\ref{th:micro}. For the rest choose a
large fixed $A$ and use consecutive half-open bins of length $s=|h|$.
Partition $(r+1)k$ such bins starting at $T$ into $r+1$ blocks of $k$
bins, where $k=\lfloor A/(s\ell)\rfloor$ and $\ell=\log\log T$.
Write $a_n$ for the integer count in bin $n$. The same counting formula
gives a block mean $\nu+O_{\chi,K_0}(1/A)+o(1)$: the endpoint error
$O_\chi(\log T/\ell)$ is divided by $k\asymp A\log T/(\nu\ell)$.
The constant in the endpoint estimate is independent of fixed $A$;
the smooth-phase variation over all the blocks is $o(1)$ in the mean.
Since $\nu>1/2$ and its distance to each positive integer is at least
$\delta$, take $A$ large enough that every block mean is strictly between
$\lfloor\nu\rfloor$ and $\lceil\nu\rceil$. Thus each block contains
a count below $\nu$ and one above it, whichever sign we request.

Choose one bin index per block with $a_n-\nu$ of alternating signs.
In the order-$r$ Lagrange divided difference, denominator signs also
alternate, so every summand has the same strict sign. Reversing the
choices gives the opposite sign. Subtracting the constant $\nu$ has no
effect on differences of order $r\ge1$. The consecutive-site identity
in the preceding proof now forces $\Delta^r a_n$ to take both signs.
For $h=s>0$, $\Delta^r a_n$ is the exact signed divisor count of
$F_{r,h}$ in bin $n$, because translation by $ks$ preserves the bin
boundaries (with the same half-open convention). For $h=-s$, reverse
the grid, or translate the bin indices by $r$ and apply the common
sign $(-1)^r$. All involved bins stay within $O_{r,K_0,\delta,\chi}
(1/\ell)$ of $T$, as asserted.
\end{proof}


\section{A finite example and further questions}
Take $\zeta$ and the three nonprincipal real primitive characters
$\chi_{-4},\chi_8,\chi_{-8}$ inducing the four character rows modulo
eight. The last three have conductors $4,8,8$. On $u\perp\one$ the
conductor functional has norm $\sqrt6\log2$, since the centered vector of
log conductors is $(\log2)(-2,0,1,1)$. For the finite prime set
\[
 P=\{3,5,7,11,13,19,23,29,31,37,43,47\}
\]
the accompanying exact calculation gives $\|b(u)\|_2>3\log2$ for every
balanced real unit vector. In fact the character sum at class $1$ is zero,
and the four residue rows form a Hadamard matrix. In each of the other
three classes one has
\[
  \sum_{p\in P,\,p\equiv a\ (8)}
     \frac{\lfloor\log_2 p\rfloor^2}{p} > \frac94.
\]
Here only, $\log_2 p$ denotes logarithm to base two. Orthogonality gives
the stated inequality, so these twelve primes already exceed the conductor
bias in both directions. The universal margin three used in the proof is
a convenience; any fixed positive margin suffices. This example also
covers the quotient $\zeta(s)L(s,\chi_{-8})/
\{L(s,\chi_{-4})L(s,\chi_8)\}$.

A separate floating-point diagnostic in the archive tests the sign and
normalization of the prime-frequency formula. With 67 tabulated zeta
ordinates in $[27,183]$, a centered, mean-zero test at $T=105$ gives
normalized prime coefficients within $9\times10^{-5}$ of $1$ for
$p=2,3$ after gamma and pole corrections. A second run uses exactly the binomial kernel of Lemma~\ref{lem:kernel}, with $M=16$ and $W=72$. Its 63 zeros give absolute errors below $3.7\times10^{-4}$ for both primes. These broad finite windows illustrate numerical frequency selection. The theorem's radius and the completeness
of a high zero table are analytic matters; the diagnostic supplies a check
of the Fourier and explicit-formula conventions.

Three questions remain. Does a fixed window length suffice for the full range of steps in Theorem~\ref{th:diff}?
Can the dependence of $C$ and $T_0$ on the conductors and the family
size be bounded at usable scales? Can one force an ordinate at which
exactly one of several specified functions vanishes, at a comparable
radius? Multiplicity inequalities alone allow common zeros of unequal
orders. The arguments here provide no simplicity or spacing assumption.

\begin{thebibliography}{99}\small
\bibitem{lean} L. Sernau et al., \emph{GRH for Dirichlet $L$-functions, in Lean 4},
commit \texttt{72fead102f663d5c7a407964a13dc0cf9d977207}.
\href{https://github.com/lsernau/grh-lean/tree/72fead102f663d5c7a407964a13dc0cf9d977207/grh}{Readable project},
\href{https://github.com/lsernau/grh-lean/blob/72fead102f663d5c7a407964a13dc0cf9d977207/grh/src/Statements.lean}{statements},
\href{https://github.com/lsernau/grh-lean/blob/72fead102f663d5c7a407964a13dc0cf9d977207/grh/GUIDE.md}{reading guide}.
\bibitem{checks} M. Rubinstein, \href{https://www.math.uwaterloo.ca/~mrubinst/alpoge_claude_grh/}{Alp\"oge/Claude proof of the Generalized Riemann Hypothesis}, independent verification, compact edition, 26 November 2026.
\bibitem{wiki} Wikipedia contributors, \href{https://en.wikipedia.org/wiki/Riemann_hypothesis}{Riemann hypothesis}, consulted 26 November 2026.
\bibitem{proveit} V. Reshetnikov, \href{https://github.com/VladimirReshetnikov/ProveIt/tree/b29490a150d3947b114d6c9daf3398300537eb64/Analysis/GeneralizedRiemannHypothesis}{ProveIt, Generalized Riemann Hypothesis collection}.
See the directory index, merged article, packet
\texttt{grh\_prime\_error\_balancing}, and the staged archive
\texttt{dirichlet\_prime\_frequency\_inverse.zip}.
\bibitem{short} ProveIt, \href{https://github.com/VladimirReshetnikov/ProveIt/blob/b29490a150d3947b114d6c9daf3398300537eb64/Analysis/GeneralizedRiemannHypothesis/Research/grh-methods-polynomial-density/short.tex}{Multiplicity differences in intervals of triple-logarithmic length}, section of \texttt{grh-methods-polynomial-density}, November 2026.
\bibitem{ccm} E. Carneiro, V. Chandee and M. B. Milinovich, \href{https://arxiv.org/abs/1503.00955}{A note on the zeros of zeta and $L$-functions}, Math. Z. \textbf{281} (2015), 315--332, Theorem 5 and the bound in $t$-aspect following its proof.
\bibitem{trans} K. Ramachandra and A. Sankaranarayanan, \href{https://matwbn.icm.edu.pl/ksiazki/aa/aa91/aa9145.pdf}{Notes on the Riemann zeta-function, II}, Acta Arith. \textbf{91} (1999).
\bibitem{brss} R. Balasubramanian, K. Ramachandra, A. Sankaranarayanan and K. Srinivas, \href{https://hrj.episciences.org/139}{Notes on the Riemann zeta function, III}, Hardy--Ramanujan J. \textbf{22} (1999).
\bibitem{pole} ProveIt, \href{https://github.com/VladimirReshetnikov/ProveIt/blob/b29490a150d3947b114d6c9daf3398300537eb64/Analysis/GeneralizedRiemannHypothesis/Research/prime_pole_windows/prime_pole_windows.tex}{Reading arithmetic weights from poles of a prime series}, Theorem 2 and Corollary 5, 26 November 2026.
\bibitem{euler} ProveIt, \href{https://github.com/VladimirReshetnikov/ProveIt/blob/b29490a150d3947b114d6c9daf3398300537eb64/Analysis/GeneralizedRiemannHypothesis/docs/incoming/local_euler_recovery.zip}{Identifying Dirichlet characters from a short window of zeros}, Theorems 1--2 and Corollary 6, 26 November 2026.
\bibitem{srinivas} K. Srinivas, \href{https://doi.org/10.4064/aa103-3-1}{Distinct zeros of functions in the Selberg class}, Acta Arith. \textbf{103} (2002), 201--207.
\bibitem{kp} J. Kaczorowski and A. Perelli, \href{https://doi.org/10.1353/ajm.1998.0012}{Functional independence of the singularities of a class of Dirichlet series}, Amer. J. Math. \textbf{120} (1998), 289--303.
\bibitem{perelli} A. Perelli, \href{https://www.rivmat.unipr.it/fulltext/2004-3s/pdf/04.pdf}{A survey of the Selberg class of $L$-functions, part II}, Riv. Mat. Univ. Parma (7) \textbf{3*} (2004), 83--118, especially Lemma 7.1.
\bibitem{bp} E. Bombieri and A. Perelli, \href{https://matwbn.icm.edu.pl/ksiazki/aa/aa83/aa8336.pdf}{Distinct zeros of $L$-functions}, Acta Arith. \textbf{83} (1998), 271--281.
\bibitem{banks} W. D. Banks, \href{https://arxiv.org/abs/2302.07073}{Comparing zeros of distinct Dirichlet $L$-functions}, Publ. Math. Debrecen \textbf{106} (2025), 327--335.
\bibitem{bl} W. D. Banks and K. Loftus, \href{https://arxiv.org/abs/2607.22930}{Unshared zeros of Dirichlet $L$-functions}, arXiv:2607.22930 (2026).
\bibitem{rs} Z. Rudnick and P. Sarnak, \href{https://doi.org/10.1215/S0012-7094-96-08115-6}{Zeros of principal $L$-functions and random matrix theory}, Duke Math. J. \textbf{81} (1996), 269--322, Proposition 2.1.
\bibitem{dav} H. Davenport, \emph{Multiplicative Number Theory}, 3rd ed., revised by H. L. Montgomery, Springer, 2000.
\end{thebibliography}
\end{document}
